\documentclass[10pt]{article}
\usepackage[margin=0.8in,top=0.65in,bottom=0.6in]{geometry}
\usepackage{amsmath,amssymb,amsthm}
\usepackage[T1]{fontenc}
\usepackage[utf8]{inputenc}
\usepackage{booktabs}
\usepackage{titlesec}
\usepackage{natbib}
\usepackage[hidelinks]{hyperref}

\titlespacing*{\section}{0pt}{1.1ex plus .2ex}{0.5ex}
\titlespacing*{\subsection}{0pt}{0.9ex plus .2ex}{0.3ex}
\titlespacing*{\paragraph}{0pt}{0.8ex plus .1ex}{1em}
\titleformat{\section}{\normalfont\large\bfseries}{\thesection}{0.6em}{}
\setlength{\parskip}{0.15ex}
\renewcommand{\baselinestretch}{0.95}
\setlength{\bibsep}{0pt plus 0.3ex}

\newcommand{\Xiac}{\Xi_{\alpha}}
\newcommand{\Ztwo}{\mathbb{Z}_2}

\title{Citation addendum and correction to \\
\emph{``The 2-Adic Sturmian Carry Constant in the Collatz Carry Equation''} \\
{\normalsize (5 June 2026)}}
\author{Elias De Jes\'us\\
\small Independent Researcher\\
\small ORCID: \href{https://orcid.org/0009-0007-0190-9143}{0009-0007-0190-9143}\\
\small \texttt{dejesuselias10@gmail.com}}
\date{September 2026}

\begin{document}
\maketitle\vspace{-2.5em}

\section*{Purpose}

This addendum accompanies the note \emph{The 2-Adic Sturmian Carry Constant in the Collatz Carry
Equation} \cite{dejesus2026sturmiancarry} (hereafter \textbf{the note}) without altering it. It
supplies a citation the note omitted, qualifies the one priority claim that omission made
unsustainable, and corrects a typographical slip in a proof together with a range hypothesis the
proof uses but the theorem statement did not record.

\textbf{No theorem, proof or computed value in the note changes.} The single mathematical edit is
the misprinted factor identified in Section~\ref{sec:typo}, which has no consequence for any
statement in the note. This addendum makes no new mathematical claim.

\section{The omitted citation, and what it contains}

\begin{quote}
J.~L\'opez and P.~Stoll, \emph{The 3x+1 conjugacy map over a Sturmian word}, Integers \textbf{9}
(2009), \#A13, 141--162 \cite{lopez2009conjugacy}.
\end{quote}

\noindent For every irrational slope $\alpha\in(0,1)$ with convergents $p_k/q_k$, and with
$1c_\alpha(j)=\lceil (j+1)\alpha\rceil-\lceil j\alpha\rceil$, their Theorem~1 gives, in $\Ztwo$,
\begin{equation}
\Phi(1c_\alpha) \;=\; -\tfrac13 \;-\; \sum_{j\ge 0} (-1)^{j+1}\,
\frac{2^{\,q_{j+1}+q_j-1}}{3\,(3^{p_{j+1}}-2^{q_{j+1}})(3^{p_j}-2^{q_j})}\,,
\label{eq:ls}
\end{equation}
where $\Phi$ is the Bernstein--Lagarias conjugacy map \cite{bernstein1996conjugacy}. Their
Corollary~2 gives a generalized continued fraction for $-1/\Phi(1c_\alpha)$, convergent in $\Ztwo$.
Their \S4 leaves $\Ztwo$ and studies the real function $\Phi_{\mathbb{R}}$, proving irrationality
of $\Phi_{\mathbb{R}}(m_\alpha)$ on $\mathbb{Q}\cap(\ln 2/\ln 3,\,1]$ and of a dual limit point
below $\ln 2/\ln 3$; the restriction is forced, since $\sum 2^{nq}/3^{np}$ converges in
$\mathbb{R}$ if and only if $\ln 2/\ln 3 < p/q \le 1$.

The denominators of \eqref{eq:ls} are the walls of two consecutive convergent shells, and the term
exponents are $q_{j+1}+q_j-1$. \textbf{These are the approximants and the approximation depths of
the note.}

The companion paper \cite{lopez2021periodicity} proves aperiodicity for parity words of
ones-density strictly above $\ln 2/\ln 3$, and its Theorem~1 shows that a rational divergent
trajectory would have density exactly $\ln 2/\ln 3$.

\section{Dictionary}

Signs were checked term by term against the note's own Definition~4.1, and the identity was
verified numerically modulo $2^{3000}$.

\begin{center}
\begin{tabular}{@{}lll@{}}
\toprule
& the note & L\'opez--Stoll \\
\midrule
slope & $\alpha=\log_2 3\approx 1.58496$ & $\alpha_{\mathrm{LS}}=\ln 2/\ln 3=\log_3 2\approx 0.630930$\\
relation & \multicolumn{2}{l}{$\alpha_{\mathrm{LS}} = 1/\alpha$}\\
word & $c_\alpha$, with $s_j(c_\alpha)=\lfloor j\alpha\rfloor$ & $1c_{\alpha_{\mathrm{LS}}}(j)=\lceil (j+1)\alpha_{\mathrm{LS}}\rceil-\lceil j\alpha_{\mathrm{LS}}\rceil$\\
constant & $\Xiac=\sum_{j\ge0}3^{-(j+1)}2^{\lfloor j\alpha\rfloor}$ & $\Phi(1c_{\alpha_{\mathrm{LS}}})$\\
\textbf{identity} & \multicolumn{2}{l}{$\boldsymbol{\Xiac = -\,\Phi\bigl(1c_{\ln 2/\ln 3}\bigr)}$}\\
shell & $(q_n,p_n)=(\text{length},\text{discharge})$ & $(p_k,q_k)=(\text{numerator},\text{denominator})$\\
relation & \multicolumn{2}{l}{$(q_n,p_n)=(p_k,q_k)$: their numerator is the note's word length}\\
approximant & $x_n=Q_n=-C_n/\delta_n$ & $F(p_k/q_k)=\phi(m_{p/q})/(2^{q}-3^{p})$\\
\bottomrule
\end{tabular}
\end{center}

\noindent\emph{Sign anchor.} \S9 of the note records the shell $(2,3)$ as $Q=-5$ with
$\Xi((1,2)^\infty)=C/\delta=5/1=5$; the convention above reproduces this.
\emph{A qualification on the approximants.} Their Definition~26 sets
$F(p/q)=\phi(m_{p/q})/(2^{q}-3^{p})$, which \emph{is} the periodic cycle value
$\Phi(m_{p/q}^{\infty})$. But the proof of their Lemma~27 gives
$-P_{2k+1}/Q_{2k+1}=F(p_{2k+1}/q_{2k+1})$ while
$-P_{2k}/Q_{2k}=F(p_{2k}/q_{2k})+g(p_{2k}/q_{2k})$. \textbf{Only the odd-indexed continued-fraction
convergents are periodic cycle values}; the even ones carry an additional gap term.

\medskip
\noindent\emph{Depth correspondence.} In the note's coordinates the term exponents
$q_{j+1}+q_j-1$ of \eqref{eq:ls} read $p_n+p_{n+1}-1$. Direct computation confirms this: the
partial sums of \eqref{eq:ls} have error valuations $2,4,10,26,83,148,568,1538,\dots$, i.e.\
$p_n+p_{n+1}-1$ for $n=1,2,3,\dots$ in the numbering of \S\S9--10 of the note. The note's
lower-convergent values $10,83,568$ are exactly the odd-indexed members of that sequence. Its
upper-convergent values $p_n-1$ ($7,64,484$) do \emph{not} appear there: \eqref{eq:ls} records the
valuation of its own tail, not $v_2(x_n+\Xiac)$ shell by shell.

\section{Qualified claim}

\paragraph{Original wording (\S1, second paragraph).}
\begin{quote}
``This note begins the study of the tower itself, and presents its first theorem-grade structural
result.''
\end{quote}

\paragraph{Corrected wording.}
\begin{quote}
``This note begins the study of the tower itself, and presents a theorem-grade structural result
for it. The result is not the first 2-adic theorem about this object. For every irrational slope,
L\'opez and Stoll \cite{lopez2009conjugacy} give a closed 2-adic series for the
Bernstein--Lagarias conjugacy value \cite{bernstein1996conjugacy} of the characteristic Sturmian
word, indexed by consecutive pairs of continued-fraction convergents; in the coordinates of this
note its term exponents are $p_n+p_{n+1}-1$, the lower-convergent depths of Theorem~7.1. What
follows is an independent derivation, sharper in one respect: it separates the upper and lower
convergent cases and pins each constant exactly, shell by shell.''
\end{quote}

\noindent The abstract carries the same qualification. The proof given in the note is independent
of \cite{lopez2009conjugacy} and was re-derived in full. Correspondingly, Remark~5.2, which credits
\cite{lothaire2002algebraic} for the classical word-agreement fact, should add that the 2-adic
consequence of that agreement --- the periodic values of the convergent blocks approximating the
conjugacy value of the characteristic word, with depths governed by consecutive convergent pairs
--- is Theorem~1 of \cite{lopez2009conjugacy}.

\section{Range of Theorem 7.1}

Theorem~7.1 is stated in the note without a range. Its proof, through Lemma~5.1, uses
$q_n\ge 2$ (the index range $1\le i\le q_n-1$ must be non-empty) and $q_{n+1}>q_n$. In the
numbering of \S\S9--10, where $(q_1,p_1)=(1,1)$, both hold from the shell $(2,3)$ onward, so
\textbf{Theorem~7.1 holds as proved for $n\ge 3$}. The two initial shells $(1,1)$ and $(1,2)$ have $q_n=1$ and lie outside the proof's scope. The
stated formula nevertheless returns the correct valuations there, $2$ and $1$ respectively, as
verified by direct computation; this is an observation, not an extension of the theorem. Nothing
in the note uses the two excluded shells.

\section{Correction}
\label{sec:typo}

\paragraph{Original (\S5, proof of Lemma 5.1, lower case).}
\begin{quote}
``At $j^*=q_n+q_{n+1}$ the mismatch occurs: $m_{j^*}=m_{q_{n+1}}=q_n-1$ and
$m_{j^*}+j^*|R_n| = q_n+q_n|R_n|-|R_{n+1}| > q_n$.''
\end{quote}

\paragraph{Corrected.}
\begin{quote}
``At $j^*=q_n+q_{n+1}$ the mismatch occurs: $m_{j^*}=m_{q_{n+1}}=q_n-1$ and
$m_{j^*}+j^*|R_n| = q_n+q_n|R_n|-q_n|R_{n+1}| > q_n$.''
\end{quote}

\noindent Indeed, with $m_{j^*}=q_n-1$ and $j^*=q_n+q_{n+1}$,
\[
m_{j^*}+j^*|R_n| \;=\; (q_n-1)+q_n|R_n|+q_{n+1}|R_n|
\;=\; (q_n-1)+q_n|R_n|+\bigl(1-q_n|R_{n+1}|\bigr)
\;=\; q_n+q_n|R_n|-q_n|R_{n+1}| ,
\]
using $q_{n+1}|R_n|+q_n|R_{n+1}|=1$, equation~(1) of the note. The factor $q_n$ on the last term
was dropped in the printed line.

\paragraph{Consequence: none.} The inequality $>q_n$ follows from $|R_n|>|R_{n+1}|$, which is
classical. Lemma~5.1, Theorem~7.1, Proposition~8.1 and every value in the computational ledger of
\S10 are unaffected.

% TODO(Elias): add pointer to irrationality note once posted
%   -- belongs at \S11 "Next targets", item (1): "an effective non-integrality statement for the
%      truncations of Xi_alpha". Leave the note's text as printed until the note is public.

\medskip\noindent\emph{Acknowledgment.} The omission was identified during an adversarial citation audit of
the Sturmian thread, assisted by Claude (Anthropic). All statements here, and any errors, are the
author's.

{\footnotesize
\bibliographystyle{plain}
\bibliography{references}
}

\end{document}
